\documentclass[11pt]{article}

\usepackage{amsmath}
\usepackage{setspace}
\usepackage{amssymb}
\usepackage{changepage}
\usepackage[a4paper, margin=0.5in]{geometry}

\setstretch{1.3}

\hbadness=10000
\setlength{\parindent}{0pt}

\newcommand\numberthis{\addtocounter{equation}{1}\tag{\theequation}}

\begin{document}
    % HEADER 
    \begin{center}
        \Large \textbf{STAT3006 Worksheet 02}
    \end{center}

    \rule{\textwidth}{0.3pt}

    % Question 1
    \large \textbf{1. Bias-Variance Decomposition}

    \begin{adjustwidth}{1em}{1em}
        % Question 1a
        (a) Let $h$ be a random estimator of a deterministic target $h*$. Assume expectations are 
        taken with respect to the randomness in $h$. Show that 

        \begin{align*}
            \mathbb{E}[(h - h^*)^2] = \text{Var}(h) + (\mathbb{E}[h] - h^*)^2
        \end{align*}

        \textbf{Answer:} We evaluate the variance of $(h - h^*)$. That is, 

        \begin{align*}
            \text{Var}(h - h^*) &= \mathbb{E}[(h - h^*)^2] - (\mathbb{E}[h - h^*])^2 \\ 
            \Rightarrow \text{Var}(h) + \text{Var}(h^*) &= \mathbb{E}[(h - h^*)^2] - 
            (\mathbb{E}[h - h^*])^2 \\ 
            \Rightarrow \text{Var}(h) &= \mathbb{E}[(h - h^*)^2] - (\mathbb{E}[h - h^*])^2 \\ 
            \Rightarrow \therefore \mathbb{E}[(h - h^*)^2] &= \text{Var}(h) + (\mathbb{E}[h - h^*])
            ^2
        \end{align*}

        % Question 1b
        (b) In the lecture, we saw that for a fixed \textit{test point} $\mathbf{x}_0$, the 
        regression tree estimator $f^{(n)}(\mathbf{x}_0)$ satisfies the conditional 
        bias-variance decomposition 

        \begin{align*}
            \mathbb{E}\left[(f^{(n)}(\mathbf{x}_0) - f^*(\mathbf{x}_0))^2\right] &= \mathbb{E}
            \bigg[\text{Var}(f^{(n)}(\mathbf{x}_0) | R^{(n)}(\mathbf{x}_0)) +
            (\mathbb{E}[f^{(n)}(\mathbf{x}_0) | R^{(n)}(\mathbf{x}_0)] - f^{*}(\mathbf{x}_0))^2
            \bigg] \numberthis
        \end{align*}

        Here $R^(n)(\mathbf{x}_0)$ denotes the random region (e.g. the leaf of the tree) containing
        $\mathbf{x}_0$, whose randomness arises from the training data. Thus (1) can be interpreted
        as the usual bias-variance decomposition applied \textit{conditionally on the random
        region}, and then averaged over the randomness of the training sample. \\

        Use the bias-variance decomposition to prove (1) \\

        \textbf{Answer:} For easier notation we set,

        \begin{align*}
            A &= f^(n)(\mathbf{x}_0) \\ 
            B &= f^*(\mathbf{x}_0) \\ 
            C &= R^{(n)}(\mathbf{x}_0)
        \end{align*}

        We want to show 

        \begin{align*}
            \mathbb{E}[(A - B)^2] = \mathbb{E}\{\text{var}(A | C)  + (\mathbb{E}[A | C] - B)^2\}
            \numberthis
        \end{align*}

        From Part (a) we have 

        \begin{align*}
            \mathbb{E}[(A - B)^2] = \text{var}(A) + (\mathbb{E}[A] - B)^2
        \end{align*}

        Using the law of total variance for $\text{var}(A)$ conditioning on $C$ we have 

        \begin{align*}
            \mathbb{E}[(A - B)^2] = \mathbb{E}[\text{var}(A | C)] + \text{var}(\mathbb{E}[A | C])
            + (\mathbb{E}[A] - B)^2 \numberthis
        \end{align*}

        Now evaluate $\text{var}(\mathbb{E}[A | C] - B)$ 

        \begin{align*}
            \text{var}(\mathbb{E}[A | C] - B) &= \mathbb{E}[(\mathbb{E}[A | C] - B)^2]
            - (\mathbb{E}\{\mathbb{E}[A | C] - B\})^2 \\ 
            \text{var}(\mathbb{E}[A | C]) &= \mathbb{E}[(\mathbb{E}[A | C] - B)^2] 
            - (\mathbb{E}[A] - B)^2 \text{ (since $B$ is fixed and using tower property)}
        \end{align*}

        Subbing this result into (3) we get 

        \begin{align*}
            \mathbb{E}[(A - B)^2] &= \mathbb{E}[\text{var}(A | C)] + \mathbb{E}[(\mathbb{E}[A | C]
            - B)^2] + (\mathbb{E}[A] - B)^2 + (\mathbb{E}[A] - B)^2 \\ 
            \mathbb{E}[(A - B)^2] &= \mathbb{E}[\text{var}(A | C)] + \mathbb{E}[(\mathbb{E}[A |
            C] - B)^2] \\ 
            \therefore \mathbb{E}[(A - B)^2] &= \mathbb{E}\bigg[\text{var}(A | C) + (\mathbb{E}[A |
            C] - B)^2\bigg]
        \end{align*}
    \end{adjustwidth}

    % Question 2
    \newpage
    \textbf{2. Entropy and Information Gain in Decision Trees} \\ 

    Decision trees are widely used classification models that recursively partition the input space.
    At each node of the tree, the algorithm chooses a split of the data that makes the resulting
    subsets as pure as possible, meaning that the observations in each subset mostly belong to a
    single class. \\ 

    To quantify the purity of a node, decision tree algorithms use an \textit{impurity} meansure.
    One of the most common meausres is \textit{entropy}, which originates from information theory.
    Intuitively, entropy measures the \textit{uncertainty} or \textit{disorder} in a set of class
    labels:

    \begin{itemize}
        \item If all observations belong to the same class, the node is perfectly pure and entropy
            is 0. 
        \item If the classes are evenly mixed, entropy is large
    \end{itemize}

    In a classification problem with $C$ classes and dataset $\mathcal{D}$, suppose we evaluate a
    split of $\mathcal{D}$ into two subsets $\mathcal{D}_l$ and $\mathcal{D}_r$ based on some feature
    and threshold. \\ 

    \begin{adjustwidth}{1em}{1em}
        % Question 2b
        (b) Show that the information gain of a split can be written as

        \begin{align*}
            \text{IG} = \text{Entropy(parent)} - \sum_{\text{child}}
            \frac{|R_{\text{child}}|}{|R_\text{parent}|} \text{Entropy(child)}
        \end{align*}

        Explain why maximising information gain is equivalent to minimizing the weighted post-split
        entropy. \\ 

        \textbf{Answer:} In decision trees, nodes can be characterised by their split point $s$ and
        class label $y$. For a given node, the information gain is given by 

        \begin{align*}
            \text{IG}(y, s) &= H(y) - H(y | s) \\ 
            \text{IG}(node) &= H(y) - H(y | s) \numberthis
        \end{align*}

        Where $H(y)$ is the entropy of $y$, and $H(y | s)$ is the conditional entropy measuring the
        remaining uncertainty in $y$ after observing $s$. Evaluating the $H(y)$ term in (4), it
        measures the entropy of the entire (remaining) dataset in the node before the split happens.
        Therefore we can write it as \\ 

        \begin{align*}
            H(y) = \text{Entropy(node)}
        \end{align*}

        For the $H(y | s)$ term, we have from the appendex that 
        
        \begin{align*}
            H(y | s) = \sum_{s'}p(s')H(y | s = s')
        \end{align*}

        Here $s'$ represents all the possible outcomes of the split $s$, i.e. which child node a
        data point falls into. $p(s')$ then represents the probability of a random observation
        from the parent lands in that child. This probability can equivalently represented as 

        \begin{align*}
            p(s') = \frac{|\mathcal{D}_\text{child}|}{|\mathcal{D}_{\text{parent}}|} =
            \frac{|R_\text{child}|}{|R_\text{parent}|}
        \end{align*}

        For $H(y | s = s')$ in (4), this is simply just the entropy of class labels in that child.
        Therefore it can equivalently represented by 

        \begin{align*}
        H(y | s = s') = \text{Entropy(child)}
        \end{align*}

        The result follows. To explain why maximising information gain is equiavelent to minimizing
        the weighted post-split entropy, we evaluate the following

        \begin{align*}
            \underset{\text{split}}{\arg\max} \text{ IG(split)} = \underset{\text{split}}{\arg\max}
            \bigg\{\text{Entropy(parent)} - \sum_{\text{child}}
            \frac{|R_{\text{child}}|}{|R_\text{parent}|}\bigg\}
        \end{align*}

        Since the entropy of the parent doesn't depend on the split, this evaluates to 

        \begin{align*}
            \underset{\text{split}}{\arg\max} \text{ IG(split)} 
                &= \underset{\text{child}}{\arg\max}\bigg\{-\sum_{\text{child}}
                \frac{|R_\text{child}|}{|R_\text{parent}|}\bigg\} \\ 
                &= \underset{\text{child}}{\arg\min}\bigg\{\sum_{\text{child}}
                \frac{|R_\text{child}|}{|R_\text{parent}|}\bigg\}
        \end{align*}
        % Question 2c 
        (c) Consider a binary classification problem with two classes: \textit{Positive} and
        \textit{Negative}. A node contains 10 observations:

        \begin{align*}
            \begin{matrix}
                \text{6 Positive}, & \text{4 Negative}
            \end{matrix}
        \end{align*}

        A candidate split produces the following subsets:

        \begin{align*}
            \begin{matrix}
                \mathcal{D}_l: \text{(4 positive, 1 negative)}, & \mathcal{D}_r: \text{(2 Positive, 3
                Negative)}
            \end{matrix}
        \end{align*}

        Compute: \\ 

        (i) the entropy of the parent node, 

        (ii) the entropy of each child node

        (iii) the conditional entropy after the split

        (iv) the information gain of the split \\ 

        Interpret the result. \\ 

        \textbf{Answer:} The entropy of the parent node is given by 

        \begin{align*}
            \text{Entropy(parent)} &= -\sum_{c = 1}^C p_c \log p_c \\ 
                                   &= -\frac{6}{10} \log\left( \frac{6}{10}\right) 
                                   - \frac{4}{10} \log \left(\frac{4}{10}\right) \\ 
                                   &\approx 0.6730
        \end{align*}

        The entropy of each child node is given by

        \begin{align*}
            \text{Entropy(left)} &= -\frac{4}{5}\log \left(\frac{4}{5}\right) - 
                \frac{1}{5} \log \left(\frac{1}{5}\right) \\ 
                                 &\approx 0.5004
                                 \\ 
            \text{Entropy(right)} &= -\frac{2}{5}\log \left(\frac{2}{5}\right) - \frac{3}{5}\log  
                \left(\frac{3}{5}\right) \\ 
                                  &\approx 0.6730
        \end{align*}

        The conditional entropy after the split is given by

        \begin{align*}
            H(y | s) &= \sum_{\text{child}} \frac{|R_\text{child}|}{|R_\text{parent}|}
            \text{Entropy(child)} \\ 
                     &= \frac{5}{10} \cdot 0.5004 + \frac{5}{10} \cdot 0.6730 \\ 
                     &\approx 0.5867
        \end{align*}

        The information gain of the split is then given by 

        \begin{align*}
            \text{IG} &= \text{Entropy(parent)} - \text{Conditional Entropy(children)} \\ 
                      &= 0.6730 - 0.5867 \\ 
                      &= 0.0863
        \end{align*}

        Interpreting this result, not much information was gained by performing the split. That is,
        we did not reduce a significant amount of uncertainty about the class labels after
        performing the split.
    \end{adjustwidth}
\end{document}
